\chapter{App H --- Requirements Detail}

\section{App H --- Requirements Detail}

This appendix expands the requirements in Chapter with the full sub-requirement specifications for each functional requirement (RF), technological requirement (TR), and non-functional requirement (NFR).

\subsection{Functional Requirements: Sub-Specifications}

\subsubsection{RF1: Design and Implement EVE-NG Topologies}

\begin{itemize}
\tightlist
\item
  \textbf{RF1.2}: Each topology must include:
\item
  Minimum 5--8 virtual machines per lab
\item
  Realistic network architecture (DMZ, internal networks, monitoring zones)
\item
  Vulnerable services configured per OWASP and NIST standards
\item
  Network segmentation with VLANs where appropriate
\item
  \textbf{RF1.3}: EVE-NG compatibility requirements:
\item
  .unl file format compatible with EVE-NG Community Edition
\item
  Additional compatibility with EVE-NG Professional Edition
\item
  KVM hypervisor support (Ubuntu 20.04+ LTS)
\item
  Minimum host specification: 16 GB RAM, 200 GB storage
\item
  \textbf{RF1.4}: Topology documentation:
\item
  Network diagram for each lab (logical and physical views)
\item
  IP addressing scheme documented
\item
  Service inventory and port mappings
\item
  Attack surface identification and justification
\end{itemize}

\subsubsection{RF2: Implement Vulnerabilities and Attack Scenarios}

\begin{itemize}
\tightlist
\item
  \textbf{RF2.5}: Vulnerability implementation standards:
\item
  All vulnerabilities must be fixable by administrators
\item
  No permanently broken or unrecoverable systems
\item
  Clear reset procedure for each vulnerability state
\item
  Automated scanning tools must successfully identify vulnerabilities
\end{itemize}

\subsubsection{RF3: Create Structured Assessment Materials}

\begin{itemize}
\tightlist
\item
  \textbf{RF3.3}: Support materials:
\item
  Student lab manuals with step-by-step guidance
\item
  Teacher guides with deployment and customisation instructions
\item
  Expected outcomes and common pitfalls documentation
\item
  Troubleshooting guides for common technical issues
\item
  \textbf{RF3.4}: Validation materials:
\item
  Automated validation scripts to verify successful lab completion
\item
  Manual verification procedures for educators
\item
  Documentation of success indicators
\end{itemize}

\subsection{Technological Requirements: Sub-Specifications}

\subsubsection{TR1: Virtualisation Infrastructure}

\begin{itemize}
\tightlist
\item
  \textbf{TR1.2}: Hypervisor:
\item
  Primary: KVM (Kernel-based Virtual Machine)
\item
  Host OS: Ubuntu Server 20.04 LTS or 22.04 LTS
\item
  QEMU integration for VM management
\item
  Nested virtualisation support for advanced scenarios
\item
  \textbf{TR1.4}: Network configuration:
\item
  Layer 2 switching capabilities within EVE-NG
\item
  VLAN support for network segmentation scenarios
\item
  IPv4 and IPv6 support
\end{itemize}

\subsubsection{TR2: Security Tools and Penetration Testing Framework}

\begin{itemize}
\tightlist
\item
  \textbf{TR2.4}: Compliance and documentation:
\item
  All tools must be freely available or have educational licences
\item
  Open-source preference for long-term sustainability
\item
  Tool version specifications documented
\item
  Legal and ethical use guidelines included
\end{itemize}

\subsubsection{TR3: Scripting and Automation}

\begin{itemize}
\tightlist
\item
  \textbf{TR3.3}: Integration points:
\item
  EVE-NG API integration for programmatic lab management
\item
  Webhook support for event-driven automation
\item
  Container orchestration (Docker) for microservices
\item
  Version control integration (Git, GitHub)
\end{itemize}

\subsection{Non-Functional Requirements: Detailed Subsections}

\subsubsection{NFR1: Performance and Efficiency}

\textbf{NFR1.1 -- Deployment Performance}: - Lab deployment time ≤ 2 minutes from topology load to ready state - VM boot time ≤ 90 seconds per VM (95\% - Reset operation ≤ 30 seconds to restore to initial state - Concurrent lab instances: support minimum 5 simultaneous labs

\textbf{NFR1.2 -- Resource Optimisation}: - VM image sizes ≤ 5 GB per image (compressed) - Memory footprint ≤ 6 GB per lab instance during execution - Minimal external dependencies (all local simulation)

\textbf{NFR1.3 -- Scalability}: - Current infrastructure supports 3--4 concurrent instances on the project HomeLab server - Architecture horizontally scalable: additional EVE-NG nodes can be added - EVE-NG Professional Edition supports larger concurrent deployments

\subsubsection{NFR2: Compliance with Security Frameworks}

\textbf{NFR2.1 -- NIST Cybersecurity Framework} : Core Functions (Identify, Protect, Detect, Respond, Recover) mapped to lab exercises; assessment procedures aligned with NIST guidelines.

\textbf{NFR2.2 -- CIS Critical Controls} : Controls 1--18 incorporated where pedagogically appropriate; assessment rubrics aligned with CIS benchmarks; remediation procedures follow CIS recommendations.

\textbf{NFR2.3 -- OWASP Standards} : OWASP Top 10 fully covered; OWASP Testing Guide methodology integrated; development practices reflected in lab design.

\textbf{NFR2.4 -- Penetration Testing Standards} : PTES phases mapped to labs; CEH curriculum alignment; ethical hacking principles strictly enforced.

\subsubsection{NFR3: Reliability and Availability}

\textbf{NFR3.1 -- System Reliability}: - MTBF ≥ 100 hours continuous operation - MTTR ≤ 5 minutes - Scenario reproducibility: 100\% - Zero critical bugs at release

\textbf{NFR3.2 -- Data Integrity}: No unintended data loss during lab operations; snapshot and restore functionality; backup procedures for sensitive configurations.

\textbf{NFR3.3 -- Availability}: Planned maintenance windows weekends only; system availability ≥ 99\% instructional periods; graceful degradation if resources become limited.

\subsubsection{NFR4: Security and Ethical Compliance}

\textbf{NFR4.1 -- Educational Sandbox Isolation}: Complete network isolation from production systems; no outbound internet access from lab environments; contained attack scenarios (no external targets).

\textbf{NFR4.2 -- Access Control and Accountability}: Role-based access control (RBAC) for lab management; audit logs of all student actions where applicable; clear identification of authorised vs.~unauthorised activities.

\textbf{NFR4.3 -- Ethical Guardrails}: Mandatory ethics training before lab access; written agreements regarding responsible use; monitoring and escalation procedures for suspicious activities.

\subsubsection{NFR5: Usability and Support}

\textbf{NFR5.1 -- Ease of Use}: Average faculty deploy time 5--10 minutes; intuitive web interface with clear navigation; minimal technical prerequisites for students.

\textbf{NFR5.2 -- Documentation and Support}: Comprehensive user manuals and quick-start guides; FAQ documentation; responsive support channel.

\textbf{NFR5.3 -- Accessibility}: WCAG 2.1 Level AA compliance for web interfaces; screen reader compatibility; keyboard navigation support; clear contrast ratios.

\subsubsection{NFR6: Institutional Sustainability}

\textbf{NFR6.1 -- Long-term Maintainability}: All components open-source or freely available; no vendor lock-in; source code in GitHub with clear documentation; backward compatibility with older EVE-NG versions where feasible.

\textbf{NFR6.2 -- Curriculum Alignment}: Explicit mapping to GEISI degree learning outcomes; integration with existing course structures; flexible customisation for different instructional contexts.

\textbf{NFR6.3 -- Scalability for Institutional Adoption}: Deployable on institutional infrastructure (not SaaS-dependent); support for multiple faculty deployments simultaneously; resource reservation and scheduling capabilities.
